\chapter{Prompt engineering}
\label{ch:prompt-engineering}

\begin{quote}
\emph{``Prompt engineering is the art of communicating with a language model. The model's capabilities are fixed; your prompt determines how much of that capability you unlock.''}
\end{quote}

% ============================================================
\section{Why prompts matter}

A large language model is a general-purpose text predictor. The prompt frames the task, provides context, and constrains the output format. The same model can be a translator, a code generator, or a medical assistant --- depending entirely on the prompt.

% ============================================================
\section{Setup: free API access}

We use the Groq free API (fast inference on Llama 3 and Mixtral) and Google Gemini free tier:

\begin{minted}{python}
# Install clients
# !pip install groq google-generativeai

import os
# Set your API keys (get free keys from console.groq.com and aistudio.google.com)
os.environ["GROQ_API_KEY"] = "your-groq-key"
os.environ["GOOGLE_API_KEY"] = "your-google-key"
\end{minted}

\begin{minted}{python}
from groq import Groq

client = Groq()

def ask_groq(prompt, model="llama-3.1-8b-instant", temperature=0.7):
    """Send a prompt to Groq and return the response."""
    response = client.chat.completions.create(
        model=model,
        messages=[{"role": "user", "content": prompt}],
        temperature=temperature,
        max_tokens=1024,
    )
    return response.choices[0].message.content

print(ask_groq("What is the capital of Benin?"))
\end{minted}

\begin{minted}{python}
import google.generativeai as genai

genai.configure()
gemini = genai.GenerativeModel("gemini-2.0-flash")

def ask_gemini(prompt, temperature=0.7):
    """Send a prompt to Gemini and return the response."""
    response = gemini.generate_content(
        prompt,
        generation_config=genai.GenerationConfig(temperature=temperature),
    )
    return response.text

print(ask_gemini("What is the capital of Benin?"))
\end{minted}

% ============================================================
\section{Zero-shot prompting}

Give the model a task with no examples:

\begin{minted}{python}
prompt = """Classify the following movie review as POSITIVE or NEGATIVE.

Review: "The cinematography was breathtaking, but the plot was
completely incoherent and the acting felt wooden."

Classification:"""

print(ask_groq(prompt, temperature=0))
\end{minted}

% ============================================================
\section{Few-shot prompting}

Provide examples to guide the model's behavior:

\begin{minted}{python}
prompt = """Translate English to French.

English: The weather is beautiful today.
French: Le temps est magnifique aujourd'hui.

English: I would like a cup of coffee, please.
French: Je voudrais une tasse de cafe, s'il vous plait.

English: Where is the nearest hospital?
French:"""

print(ask_groq(prompt, temperature=0.2))
\end{minted}

\begin{aitip}
Few-shot examples should be diverse, high-quality, and representative of the expected input. 3--5 examples is usually sufficient. More examples can improve accuracy on structured tasks.
\end{aitip}

% ============================================================
\section{Chain-of-thought (CoT) prompting}

Ask the model to reason step by step before giving a final answer:

\begin{minted}{python}
# Without CoT
prompt_direct = "If a train travels 120 km in 1.5 hours, and then 80 km in 1 hour, what is the average speed for the entire trip?"

# With CoT
prompt_cot = """If a train travels 120 km in 1.5 hours, and then 80 km in 1 hour, what is the average speed for the entire trip?

Let's think step by step:"""

print("=== Direct ===")
print(ask_groq(prompt_direct, temperature=0))
print("\n=== Chain-of-Thought ===")
print(ask_groq(prompt_cot, temperature=0))
\end{minted}

Chain-of-thought improves accuracy on math, logic, and multi-step reasoning tasks dramatically.

% ============================================================
\section{Role prompting}

Assign a role to the model using the system message:

\begin{minted}{python}
from groq import Groq
client = Groq()

response = client.chat.completions.create(
    model="llama-3.1-8b-instant",
    messages=[
        {"role": "system", "content": (
            "You are an expert data scientist. You explain concepts "
            "clearly using analogies. You always provide Python code "
            "examples. You never use jargon without defining it first."
        )},
        {"role": "user", "content": "Explain overfitting to a beginner."},
    ],
    temperature=0.7,
)
print(response.choices[0].message.content)
\end{minted}

% ============================================================
\section{Structured output}

Force the model to produce JSON, CSV, or other structured formats:

\begin{minted}{python}
prompt = """Extract the following information from the text and return
it as a JSON object with keys: name, age, condition, medication.

Text: "Mrs. Fatou Diallo, 67 years old, was diagnosed with
Type 2 diabetes last year. She takes Metformin 500mg twice daily."

JSON:"""

import json
result = ask_groq(prompt, temperature=0)
print(result)
data = json.loads(result)
print(f"Patient: {data['name']}, Age: {data['age']}")
\end{minted}

\begin{minted}{python}
# Using Gemini with structured output config
import google.generativeai as genai
import typing_extensions as typing

class PatientInfo(typing.TypedDict):
    name: str
    age: int
    condition: str
    medication: str

model = genai.GenerativeModel("gemini-2.0-flash")
result = model.generate_content(
    "Extract patient info: Mrs. Fatou Diallo, 67, Type 2 diabetes, Metformin 500mg",
    generation_config=genai.GenerationConfig(
        response_mime_type="application/json",
        response_schema=PatientInfo,
    ),
)
print(result.text)
\end{minted}

% ============================================================
\section{Prompt templates}

Build reusable prompts with variables:

\begin{minted}{python}
from string import Template

summarizer = Template("""Summarize the following $doc_type in $num_sentences sentences.
Focus on: $focus_area

Text:
$text

Summary:""")

prompt = summarizer.substitute(
    doc_type="research paper abstract",
    num_sentences="3",
    focus_area="methodology and key findings",
    text="We propose LoRA, a method for adapting large language models..."
)
print(ask_groq(prompt))
\end{minted}

\begin{minted}{python}
# With LangChain prompt templates
from langchain_core.prompts import ChatPromptTemplate

template = ChatPromptTemplate.from_messages([
    ("system", "You are a {role}. Respond in {language}."),
    ("human", "{question}"),
])

prompt = template.invoke({
    "role": "medical doctor",
    "language": "French",
    "question": "What are the symptoms of malaria?",
})
print(prompt.to_string())
\end{minted}

% ============================================================
\section{Common prompt patterns}

\begin{longtable}{p{3.5cm}p{9cm}}
\toprule
\textbf{Pattern} & \textbf{Example} \\
\midrule
Persona & ``You are a senior Python developer...'' \\
Task framing & ``Your task is to classify emails as spam or not-spam.'' \\
Output format & ``Return your answer as a JSON array.'' \\
Constraints & ``Use only information from the provided text.'' \\
Step-by-step & ``Think step by step before answering.'' \\
Self-consistency & Run 3 CoT paths, take the majority answer. \\
\bottomrule
\end{longtable}

\begin{exercise}
\begin{enumerate}
  \item Write a zero-shot prompt that classifies customer support tickets into categories: billing, technical, account, general. Test on 5 example tickets.
  \item Create a few-shot prompt (3 examples) for named entity recognition. The model should extract person names, organizations, and locations from a text.
  \item Use chain-of-thought prompting to solve: ``A store sells apples at \$1.50 each. A customer buys 3 apples and pays with a \$10 bill. They also have a 20\% discount coupon. How much change do they receive?''
  \item Build a reusable prompt template for code review. The template should accept: language, code snippet, and review focus (security, performance, readability).
  \item Compare the same prompt on Groq (Llama 3) and Gemini. Note differences in format, verbosity, and accuracy.
\end{enumerate}
\end{exercise}

\begin{ethicsbox}
Prompt engineering can be used to bypass safety guardrails (``jailbreaking''). As AI practitioners, we have a responsibility not to develop or share jailbreak prompts. When testing model safety, use controlled environments and report vulnerabilities to model providers through responsible disclosure.
\end{ethicsbox}

% ============================================================
\section{Chapter summary}

\begin{itemize}
  \item The prompt is the interface between human intent and model behavior.
  \item Zero-shot works for simple tasks; few-shot adds examples for harder tasks.
  \item Chain-of-thought prompting dramatically improves reasoning by asking the model to show its work.
  \item Role prompting via system messages controls tone, expertise, and format.
  \item Structured output (JSON, schemas) makes LLM responses machine-readable.
  \item Prompt templates make prompts reusable and parameterizable.
  \item Free APIs (Groq, Gemini) provide access to state-of-the-art models at no cost.
\end{itemize}
